\scanpage{28}
§ 4 \uline{Site de De Rham}

\diagram{\begin{tikzcd}[ampersand replacement=\&]
 \& U \arrow[r,phantom,"\subset" description] \& X \arrow[dd,"f"] \\
U' \arrow[ur] \& \& \\
 \& \& S
\end{tikzcd}}

Schéma relatif. Les objets du \uline{site de De Rham} \inserted{$(X/S)_{\mathrm{DR}}$, ou $\mathrm{DR}(X/S)$} sont les \deleted{triplets} \inserted{quadruples} \deleted{$(U,\mathcal A,\varphi,\pi^*)$}\inserted{$(U,U',i,\pi)$}, où $U$ est un ouvert de $X$, \deleted{$\mathcal A$ un faisceau \unsure{de} $i:U\longrightarrow U'$} $U'$ un $S$-préschéma, $i:U\longrightarrow U'$ une $S$-imm. nilpotente, et $\pi$ une structure de puissances divisées tronquées (pour $n$ convenable) sur $J=\operatorname{Ker}(\underline{O}_{U'}\unsure{-|\longrightarrow}\underline{O}_U)$. Morphismes du site évidents. Topologie itou: la ``topologie de Zariski''.

Le site de \uline{De Rham restreint} (ou \inserted{\uline{site de De Rham}} \uline{\unsure{stratifiant|strictifiant}}) est formé des objets tels que localement sur $U'$ il \uline{existe} une rétraction sur $U$ (N.B. cette rétraction n'entre pas dans la structure).

Topos. $\widetilde{(X/S)}_{\mathrm{DR}}$; $\widetilde{(X/S)}_{\mathrm{DRst}}$ ou \deleted{simplement} $\widetilde{\mathrm{DR}}(X/S)$, $\widetilde{\mathrm{DRst}}(X/S)$; si une ambiguïté n'est pas à craindre, on laisse tomber le $/S$.

Diagramme de morphismes de topos
\mspage{27}
